\section{Perspectiva del transporte óptimo}

La causa fundamental de la curvatura de las trayectorias radica en la \textbf{forma de emparejamiento} de $\mathbf{x}_0$ y $\mathbf{x}_1$. En esta sección se aborda este problema desde la perspectiva del transporte óptimo (Optimal Transport).

\subsection{Emparejamiento aleatorio vs.\ emparejamiento óptimo}

En el entrenamiento estándar de Flow Matching, $\mathbf{x}_0$ y $\mathbf{x}_1$ se muestrean de forma \textbf{independiente}:
$$\mathbf{x}_0 \sim p_0, \quad \mathbf{x}_1 \sim p_1, \quad \text{ambos independientes}$$

Esto significa que una muestra ruidosa $\mathbf{x}_0$ podría ser ``asignada'' a cualquier muestra de datos $\mathbf{x}_1$, sin preferencia alguna por la estructura espacial.

\begin{importantbox}{Idea fundamental del transporte óptimo (Optimal Transport)}
El problema del transporte óptimo se centra en: ¿cómo encontrar el plan de transporte más ``económico'' desde la distribución $p_0$ hasta la distribución $p_1$? Específicamente, para un costo $L_2$ (también conocido como distancia Wasserstein-2), el plan de transporte óptimo minimiza:
$$W_2^2(p_0, p_1) = \inf_{\pi \in \Pi(p_0, p_1)} \int \|\mathbf{x}_0 - \mathbf{x}_1\|^2 \, d\pi(\mathbf{x}_0, \mathbf{x}_1)$$
donde $\Pi(p_0, p_1)$ es el conjunto de todos los acoplamientos (distribuciones conjuntas), cuyas márgenes son $p_0$ y $p_1$, respectivamente.

\medskip
\textbf{Intuición}: El transporte óptimo tiende a emparejar puntos ``cercanos'', evitando desplazamientos de larga distancia, lo que minimiza así el costo total del traslado.
\end{importantbox}

\subsection{Influencia del emparejamiento OT en la rectitud de las trayectorias}

Si utilizamos un plan de transporte óptimo para emparejar $\mathbf{x}_0$ y $\mathbf{x}_1$, en lugar de un emparejamiento aleatorio, entonces:

\begin{itemize}
    \item El emparejamiento es más ``racional'': los puntos ruidosos cercanos tienden a mapearse hacia puntos de datos cercanos.
    \item Hay menos cruces entre las trayectorias condicionales de distintos emparejamientos.
    \item Las colisiones en los campos vectoriales marginales son menores, lo que hace que las trayectorias marginales sean naturalmente más rectas.
\end{itemize}

\begin{warningbox}{Costo computacional del emparejamiento OT}
Aunque el emparejamiento OT es teóricamente muy atractivo, la complejidad de calcular exactamente un plan de transporte óptimo es muy alta. Para $n$ muestras, la complejidad de los algoritmos clásicos es $O(n^3)$, lo cual no es práctico en entrenamientos a gran escala. Por lo tanto, en la práctica se suelen emplear métodos aproximados, como calcular OT dentro de un mini-lote (Mini-batch OT), o utilizar Reflow como alternativa; Reflow puede considerarse una optimización de emparejamiento OT ``implícita''.
\end{warningbox}

\subsection{Relación entre Reflow y OT}

Existe una conexión profunda entre Reflow y el transporte óptimo:

\begin{itemize}
    \item Reflow optimiza implícitamente el plan de emparejamiento $\mathbf{x}_0 \to \mathbf{x}_1$ mediante ajuste fino iterativo.
    \item Cada ronda de Reflow acerca el emparejamiento al plan OT, ya que el mapeo aprendido por el modelo tiende a seguir la ``ruta más corta''.
    \item Teóricamente, el límite de infinitas iteraciones de Reflow es el mapeo OT (bajo ciertas condiciones de regularidad).
\end{itemize}

\begin{knowledgebox}{Curvatura de las trayectorias vista desde el costo de transporte}
La esencia de la curvatura de las trayectorias radica en un \textbf{costo de transporte excesivo}. Cuando se emparejan puntos distantes, la línea que los conecta atraviesa regiones intermedias, cruzándose con otras líneas rectas de distintos emparejamientos. Estos cruces generan inconsistencias en el campo vectorial marginal (múltiples direcciones de velocidad en el mismo punto); la red solo puede aprender la velocidad esperada (promedio), lo que resulta en trayectorias curvas. El emparejamiento OT reduce los cruces al minimizar el costo de transporte, mientras que Reflow logra un efecto similar mediante una optimización implícita iterativa.
\end{knowledgebox}

\subsection{Resumen de la sección}

El transporte óptimo proporciona un marco teórico para comprender el problema de la curvatura de las trayectorias. El emparejamiento aleatorio conduce a un alto costo de transporte y cruces en las trayectorias; el emparejamiento OT puede aliviar este problema, aunque su costo computacional es elevado. Reflow ofrece una alternativa práctica que optimiza implícitamente el emparejamiento mediante ajuste fino iterativo.

````
